\chapter{Implementation}\label{ch:ch06}

The previous chapter fixed the shape of the system: four request layers, three
authorization gates, fifteen entities, and a single Gin process fronting
PostgreSQL and an external Judge0 sandbox. This chapter descends from that
blueprint to the code that realizes it. It follows the two halves of APEX in
turn --- the Go backend and the React frontend --- and traces the flows that
carry the most engineering weight: the teacher's exam builder, the timed
exam-taking session with server-side autosave and resume, the exam lifecycle
and attempt state machine, the notification fan-out, and the in-process
reminder scheduler. The auto-grading pipeline is deliberately \emph{not} covered
here; it is substantial enough to warrant its own treatment in
\cref{ch:ch07}, and this chapter references it only where a flow hands off to
it.

The figures throughout are drawn from the real code paths. Wherever a number,
identifier, route, or column name appears, it is quoted verbatim from the
repository rather than paraphrased.

\section{Backend Handler Anatomy}

Every backend package that owns an HTTP surface follows the same internal
shape. A package such as \texttt{internal/exam} exposes a \texttt{RegisterRoutes}
function that mounts its endpoints onto the shared \texttt{public} and
\texttt{protected} route groups, and a set of handler functions that Gin
invokes per request. The route file is deliberately thin --- it is the single
place where a reader can see the whole feature's URL surface and its role gate
at a glance:

\begin{lstlisting}[language=Go,caption={The exam feature's route registration --- one place, one role gate.},label={lst:examroutes-reg}]
func RegisterRoutes(public, protected *gin.RouterGroup) {
    exams := protected.Group("/exams")
    exams.Use(middleware.RequireRole("teacher"))
    {
        exams.POST("", CreateExam)
        exams.GET("", GetExams)
        exams.GET("/:id", GetExam)
        exams.PUT("/:id", UpdateExam)
        exams.DELETE("/:id", DeleteExam)
        exams.POST("/:id/assign", AssignExam)
        exams.POST("/:id/close", CloseExam)
        exams.POST("/:id/reopen", ReopenExam)
        exams.GET("/:id/results", GetExamResults)
    }
}
\end{lstlisting}

The group is already on the authenticated \texttt{protected} chain, so
\texttt{middleware.RequireRole("teacher")} is the only additional gate the
package declares. Gin~\cite{gin} is the HTTP framework throughout; its router
group mechanism is what lets the two cross-cutting concerns --- authentication
and role --- be expressed once and inherited by every route below.

Inside the handlers, a consistent four-step spine recurs, mapping directly onto
the request lifecycle established in the design chapter:

\begin{enumerate}
  \item \textbf{Bind and validate.} The handler binds the JSON body into a
    request struct with \texttt{c.ShouldBindJSON}, relying on Gin's
    \texttt{binding:"required"} tags for the mandatory fields and returning
    \texttt{400} with a terse, non-leaking message otherwise.
  \item \textbf{Identify the caller.} \texttt{c.GetUint("user\_id")} reads the
    subject that the auth middleware placed on the context; the role is read the
    same way when a handler serves both teacher and student.
  \item \textbf{Enforce resource ownership.} The handler re-queries the target
    row scoped to the caller --- almost always as
    \texttt{Where("id = ? AND teacher\_id = ?", id, teacherID)} --- so that a
    missing row and a row owned by somebody else are indistinguishable and both
    surface as \texttt{404}. This is the third authorization gate, and it lives
    in the handler because only the handler knows what ``owns'' means for its
    resource.
  \item \textbf{Mutate, transactionally when needed.} Single-row writes go
    straight through GORM~\cite{gorm}; multi-table writes are wrapped in
    \texttt{database.DB.Transaction(func(tx *gorm.DB) error \{ ... \})} so a
    partial failure rolls back cleanly.
\end{enumerate}

Two handlers show the pattern at its extremes. \texttt{GradeSubmission} (manual
grading of a written answer) is a small handler whose entire body is an
ownership proof: it loads the submission, loads the exam it belongs to, and
refuses with \texttt{403} unless \texttt{exam.TeacherID} equals the caller.
Only then does it write the score and re-run attempt finalization:

\begin{lstlisting}[language=Go,caption={Ownership before mutation: a teacher may grade only submissions on exams they own.},label={lst:grade-ownership}]
// Verify the teacher owns the exam this submission belongs to.
var exam models.Exam
if err := database.DB.Select("id", "teacher_id").First(&exam, sub.ExamID).Error; err != nil {
    c.JSON(http.StatusNotFound, gin.H{"error": "exam not found"})
    return
}
if exam.TeacherID != teacherID {
    c.JSON(http.StatusForbidden, gin.H{"error": "not authorized to grade this submission"})
    return
}
sub.Score = req.Score
sub.Status = req.Status
// ... persist, then hand off to the grading aggregator:
if sub.ExamAttemptID != nil {
    judge0.FinalizeAttempt(*sub.ExamAttemptID)
}
\end{lstlisting}

At the other end, \texttt{DeleteExam} is the largest transactional handler. An
exam is the root of a small object graph --- problems, test cases, attempts,
submissions, test results, and class assignments --- and PostgreSQL foreign
keys are only partially declared, so the handler deletes children explicitly in
dependency order inside one transaction: test results (selected through their
parent submissions), test cases, submissions, attempts, problems, exam--class
links, and finally the exam row. If any step errors, the whole delete rolls
back and the client receives \texttt{500} with the underlying message; the
half-deleted state never becomes visible.

A subtler variant of ownership appears in \texttt{AssignExam}, which links an
exam to one or more classes. A naive implementation would trust the
\texttt{class\_ids} array from the request body. Instead the handler counts how
many of those IDs actually belong to the calling teacher and rejects the whole
request with \texttt{403} unless every one checks out:

\begin{lstlisting}[language=Go,caption={Batch ownership check in AssignExam --- all class IDs must belong to the caller.},label={lst:assign}]
var ownedCount int64
database.DB.Model(&models.Class{}).
    Where("id IN ? AND teacher_id = ?", req.ClassIDs, teacherID).
    Count(&ownedCount)
if int(ownedCount) != len(req.ClassIDs) {
    c.JSON(http.StatusForbidden, gin.H{"error": "one or more classes do not belong to you"})
    return
}
\end{lstlisting}

This closes an insecure-direct-object-reference class of bug at the point where
it would otherwise open: a teacher cannot push their exam into a colleague's
class by guessing a class ID.

\section{The Exam Builder}

The exam builder (\texttt{frontend/src/features/exams/pages/ExamBuilder.tsx})
is the most stateful screen in the product and the one that exercises the most
backend endpoints in a single user action. It is a client-orchestrated
composition: rather than a single ``save exam'' call, the builder drives the
exam, problem, test-case, and assignment endpoints in sequence, treating the
backend as a set of primitives and owning the orchestration itself.

The screen holds the entire draft exam in React state --- title, description,
duration, passing score, the shuffle and show-results toggles, the schedule
(start and close date/time), the assigned course, the draft flag, and an
ordered array of \texttt{Question} objects. New questions are created locally by
\texttt{createQuestion(type)} and only materialize on the server at save time.
The builder also imports from the teacher's shared question bank: bank problems
are filtered client-side by a search string and a type filter, selected, and
mapped into the same \texttt{Question} shape via \texttt{bankProblemToQuestion}.

Validation runs before any network call. \texttt{handleSaveClick} refuses to
proceed unless there is a non-empty title, a valid assigned course, at least one
question, and --- critically --- every MCQ has at least one correct option
selected. The last check is the frontend mirror of a backend guarantee from the
grading chapter: an MCQ with no configured answer can only ever land in
\texttt{pending\_review}, so catching it at authoring time spares the teacher a
queue of un-gradable questions.

When validation passes, \texttt{handleSave} executes a three-phase orchestration:

\begin{enumerate}
  \item \textbf{Upsert the exam.} On create it \texttt{POST}s to \texttt{/exams};
    on edit it \texttt{PUT}s to \texttt{/exams/:id} and then deletes every
    existing problem (a deliberately simple ``replace'' strategy --- the builder
    re-creates the full problem set rather than diffing it).
  \item \textbf{Create each problem.} For every question it \texttt{POST}s to
    \texttt{/exams/:id/problems}, projecting the union-typed \texttt{Question}
    down to the flat problem payload the backend expects. MCQ questions carry
    JSON-encoded \texttt{options} and \texttt{correct\_option\_ids}; written
    questions carry \texttt{max\_word\_count}, \texttt{rubric}, and
    \texttt{require\_manual\_grading}; coding questions carry
    \texttt{starter\_code}, \texttt{time\_limit\_ms}, and
    \texttt{memory\_limit\_kb}. For coding problems it then \texttt{POST}s each
    non-empty test case to \texttt{/problems/:id/test-cases}, preserving order
    via an explicit \texttt{order\_index}.
  \item \textbf{Assign to the class.} Finally it \texttt{POST}s to
    \texttt{/exams/:id/assign} with the single selected class ID.
\end{enumerate}

The schedule is assembled from separate date and time-of-day inputs into an ISO
timestamp; if the teacher sets a start but no close time, the builder derives the
close as start plus the exam's duration. The draft flag defaults to the
teacher's stored preference (\texttt{default\_exam\_draft}, itself defaulting to
\texttt{true}), so exams are born unpublished unless the teacher opts otherwise.
Because \texttt{shuffle\_questions} is persisted on the exam row, it travels
with the exam through create, edit, and re-hydration on the edit screen; the
canonical stored order remains \texttt{order\_index}, which every read path
preloads in ascending order.

One further subtlety guards data integrity. If the teacher edits an exam that
students have already attempted (\texttt{attempt\_count > 0}), the builder does
not silently save. It opens a confirmation dialog, because a structural edit is
destructive: it triggers the server's reset mechanism, which wipes attempts,
submissions, and test results and stamps \texttt{exams.reset\_at}. Surfacing
that consequence at the moment of editing is a small piece of honesty that
prevents a teacher from erasing a class's work by accident.

\section{Exam Lifecycle and Attempt State Machine}

Two related state machines govern an exam's life, shown together in
\cref{fig:exam-state}. The upper machine is the \emph{exam} lifecycle as seen by
the teacher; the lower one is the \emph{attempt} state machine as experienced by
one student's sitting.

\begin{figure}[htbp]
\centering
\resizebox{\textwidth}{!}{%
\begin{tikzpicture}[
    node distance=16mm and 20mm,
    every node/.style={font=\sffamily\small},
    guard/.style={apexlabel},
    trans/.style={font=\sffamily\scriptsize, midway, above, sloped},
    transb/.style={font=\sffamily\scriptsize, midway, below, sloped},
  ]

  % ---- Exam lifecycle (top row) ----
  \node[apexlabel, anchor=west] (lblexam) at (0,1.35) {Exam lifecycle};
  \node[apexbox] (draft) {Draft\\\scriptsize\itshape is\_draft = true};
  \node[apexstore, right=of draft] (published) {Published\\\scriptsize\itshape is\_draft = false};
  \node[apexext, right=of published] (reset) {Reset\\\scriptsize\itshape reset\_at set};

  \draw[apexarrowblue] (draft) -- node[trans] {publish}
        node[transb] {\itshape [start\_time $\neq$ nil]} (published);
  \draw[apexarrow] (published.north east) to[out=35,in=145]
        node[trans] {teacher destructive edit} (reset.north west);
  \draw[apexarrow, dashed] (reset.south west) to[out=-145,in=-35]
        node[transb] {wipe attempts, re-open} (published.south east);

  % ---- Attempt state machine (bottom row) ----
  \node[apexlabel, anchor=west] (lblatt) at (0,-2.65) {Attempt state machine};
  \node[apexbox, below=32mm of draft] (inprog) {in\_progress};
  \node[apexbox, right=of inprog] (submitted) {submitted};
  \node[apexstore, right=of submitted] (grading) {grading / running};
  \node[apexstore, right=of grading] (graded) {graded};

  % autosave self-loop
  \draw[apexarrow] (inprog.north) to[out=120,in=60,looseness=6]
        node[trans, above=1pt] {PUT autosave (draft\_answers)} (inprog.north);

  \draw[apexarrowblue] (inprog) -- node[trans] {submit}
        node[transb] {\itshape POST submit} (submitted);
  \draw[apexarrowblue] (submitted) -- node[trans] {grade (async)} (grading);
  \draw[apexarrowblue] (grading) -- node[trans] {finalize}
        node[transb] {\itshape no pending\_review} (graded);
  \draw[apexarrow] (grading.north) to[out=110,in=70,looseness=3]
        node[trans, above=1pt] {per test case (Judge0)} (grading.north);

  % link: publishing enables attempts
  \draw[apexarrow, dashed, draw=apexgrey] (published.south) --
        node[apexlabel, midway] {student opens exam} (inprog.north east);

\end{tikzpicture}%
}
\caption{Exam lifecycle and attempt state machine.}
\label{fig:exam-state}
\end{figure}


An exam is created as a \textbf{draft} (\texttt{is\_draft = true}). Publishing
is the transition \texttt{is\_draft = false}, and \texttt{UpdateExam} refuses it
unless the exam has a \texttt{start\_time}: the handler computes whether the
exam ``will have a start'' after the update and returns
\texttt{400 "set a start time before publishing"} otherwise. This guard exists
because an earlier version silently defaulted an unset start to ``now,'' which
surprised teachers by opening an exam the instant they clicked Publish. The
derived \texttt{status} a client sees --- \texttt{draft}, \texttt{upcoming},
\texttt{active}, or \texttt{completed} --- is computed by
\texttt{computeExamStatus} purely from \texttt{is\_draft}, \texttt{start\_time},
\texttt{end\_time}, and the current time, never stored, so it can never drift
out of sync with the schedule.

Two teacher actions bend the timeline after publication. \texttt{CloseExam}
stamps \texttt{end\_time} to now, ending the window immediately.
\texttt{ReopenExam} extends \texttt{end\_time} by a requested number of minutes
and, if the exam had not yet opened, pulls \texttt{start\_time} forward to now so
the window is genuinely live. The most consequential transition is the dashed
loop in the figure: a destructive edit moves the exam through the \textbf{reset}
state --- \texttt{reset\_at} is stamped, attempts and their descendants are
wiped in a transaction, and class grade announcements are rolled back --- before
returning to a clean published state.

The attempt machine is where most of the runtime action lives. An
\texttt{ExamAttempt} row is unique per \texttt{(user\_id, exam\_id)} via
\texttt{idx\_user\_exam\_attempt}, so a student has at most one attempt per
exam. \texttt{StartAttempt} is written to be idempotent: it looks up the
existing attempt and only creates one --- with \texttt{status = "in\_progress"}
--- when none exists, after checking enrollment, assignment, publication, and
that the current time falls inside the exam window. Repeated ``start'' calls,
whether from a double-click or a cross-device resume, converge on the same row
rather than erroring or forking.

From \texttt{in\_progress} the attempt has one forward transition, \textbf{submit},
plus a self-loop for autosave (discussed next). \texttt{SubmitAttempt} is
guarded against replay: if the attempt is already \texttt{submitted} it returns
\texttt{409 "attempt already submitted"}. On a valid submit it creates one
\texttt{Submission} per answered question, dispatches each to the grading
engine by type, flips the attempt to \texttt{submitted} with a
\texttt{submitted\_at} stamp, and --- notably --- clears the autosave draft by
setting \texttt{draft\_answers} and \texttt{draft\_saved\_at} back to
\texttt{NULL} in the same update. The remaining transitions,
\texttt{grading}/\texttt{running} to \texttt{graded}, are owned entirely by the
aggregation logic of \cref{ch:ch07} and are shown here only to close the loop.

\section{The Timed Exam-Taking Flow}

\texttt{ExamTaking.tsx} is the single largest component in the frontend and the
one with the most demanding correctness requirements: it must run a countdown
that cannot be gamed by a refresh, survive a browser crash or a device switch
without losing answers, and hand a clean payload to the server exactly once. It
achieves this with a two-tier persistence strategy --- \texttt{localStorage} as
an instant offline backup, and a server-side draft as the authoritative,
cross-device source of truth.

\subsection{Server-Side Autosave}

The authoritative half of persistence is the \texttt{PUT
/student/exams/:id/autosave} endpoint, backed by two columns on the attempt
row: \texttt{draft\_answers} (a JSONB payload) and \texttt{draft\_saved\_at} (a
timestamp). The server treats the draft as \emph{opaque}: the client serializes
whatever shape its UI needs --- the answers array, the current question index,
the visited-question set, and the session start time --- and the handler stores
the raw bytes verbatim, so the payload round-trips unchanged and the server
never needs to understand the client's model. The handler's only obligations
are to locate the caller's active attempt, refuse once it has been submitted,
and validate that the body is well-formed JSON:

\begin{lstlisting}[language=Go,caption={AutosaveAttempt persists an opaque JSON draft on the active attempt.},label={lst:autosave}]
if attempt.Status == "submitted" {
    c.JSON(http.StatusConflict, gin.H{"error": "attempt already submitted"})
    return
}
// Read raw body so any client-shape JSON round-trips unchanged.
body, err := io.ReadAll(c.Request.Body)
if err != nil || len(body) == 0 {
    c.JSON(http.StatusBadRequest, gin.H{"error": "invalid body"})
    return
}
if !json.Valid(body) {
    c.JSON(http.StatusBadRequest, gin.H{"error": "body must be valid JSON"})
    return
}
now := time.Now()
database.DB.Model(&attempt).Updates(map[string]any{
    "draft_answers":  string(body),
    "draft_saved_at": now,
})
\end{lstlisting}

On the client the autosave is debounced. A \texttt{useEffect} watching the
answers, current index, and visited set schedules a \texttt{setTimeout} of
1500\,ms; each change resets the timer, so a burst of typing
produces one write rather than dozens. The call is deliberately best-effort ---
its \texttt{.catch} swallows failures --- because \texttt{localStorage} already
covers the offline case, and a transient network error during an exam must never
interrupt the student. A separate, undebounced effect writes the same snapshot
to \texttt{localStorage} on every change, so the local copy is always at least
as fresh as the last render.

\subsection{Resume and Reconciliation}

Resume is where the two tiers meet, and the reconciliation logic is the most
carefully written part of the component. On load, once the questions are known,
the component reads both snapshots --- the \texttt{localStorage} session and the
server's \texttt{draft\_answers} from the attempt included in the
\texttt{GET /student/exams/:id} response --- and chooses the newer one by
comparing timestamps (\texttt{draft\_saved\_at} for the server, the session's
\texttt{startedAt} locally). The winning snapshot's answers are merged by
\texttt{questionId} onto a freshly initialized answer array, so questions that
did not exist when the draft was saved still get clean defaults. This is what
lets a student begin an exam on a laptop, close it, and resume on a phone with
every answer intact.

Two guards protect the resume path. First, the remaining time is recomputed
from the stored \texttt{startedAt} against the exam duration --- a snapshot is
only restored if time actually remains, so a refresh cannot reset the clock.
Second, if the exam's \texttt{reset\_at} is newer than the snapshot's start
time, the teacher has destructively edited the exam since the draft was written,
and the stale cache is discarded outright rather than resurrecting answers to
questions that no longer exist. The submit path closes the loop by removing the
\texttt{localStorage} key on success, while the server-side draft is nulled by
\texttt{SubmitAttempt} itself.

The countdown timer is a one-second \texttt{setInterval} decrementing
\texttt{timeLeft}; when it reaches zero, a watching effect fires
\texttt{doSubmit(true)}, an automatic submission indistinguishable from a manual
one except for the toast it shows. \texttt{doSubmit} is idempotent via a
\texttt{submitted}/\texttt{isSubmitting} pair of flags so that the auto-submit
and a last-second manual click cannot both fire. The payload it builds skips
unanswered questions entirely: only questions with a selected option, non-empty
text, or non-empty code produce an entry in the \texttt{answers} array sent to
\texttt{POST /student/exams/:id/submit}.

\section{Question Editors and the Monaco Integration}

Each question type has a dedicated editor component, dispatched by the builder
from a discriminated union. \texttt{MCQEditor} manages an options list (capped
at six), tracks which option IDs are correct, toggles single- versus
multiple-correct, and supports an optional explanation and an inline image
(uploaded as a base64 data URL, rejected above 2\,MB).
\texttt{WrittenEditor} is the simplest --- a prompt, a maximum word count
defaulting to 500, an optional rubric, and a \texttt{require\_manual\_grading}
switch --- reflecting that written answers always enter the manual queue.
\texttt{CodingEditor} is the richest: it manages the problem statement, starter
code across the supported languages (\texttt{python}, \texttt{javascript},
\texttt{c}, \texttt{cpp}), per-problem time and memory limits, and a list of
test cases each flagged sample-or-hidden with its own point value. Sample test
cases are the ones surfaced to students for live feedback; hidden ones run only
during real grading.

The code surfaces on both sides --- the coding editor for authors and the
in-exam solution pane for students --- are the same Monaco editor,
integrated through \texttt{@monaco-editor/react}~\cite{monaco}. Monaco is the
editor that powers VS Code, and embedding it gives students syntax
highlighting, bracket matching, and familiar keybindings inside the browser
with almost no configuration. In \texttt{ExamTaking} it is mounted with a
minimal, exam-appropriate option set --- minimap disabled, word wrap on,
\texttt{automaticLayout} enabled so it reflows with the panel, and a theme
bound to the app's light/dark mode:

\begin{lstlisting}[language=TypeScript,caption={Monaco embedded in the in-exam solution pane.},label={lst:monaco}]
<Editor
  height="100%"
  language={(ans.language || "python3") === "python3" ? "python"
    : (ans.language || "python3") === "cpp" ? "cpp"
    : (ans.language || "python3")}
  value={ans.code || (q.starterCode?.[ans.language || "python3"] || "")}
  onChange={(v) => updateAnswer(currentIdx, { code: v || "" })}
  theme={theme === "dark" ? "vs-dark" : "light"}
  options={{
    fontSize: 13,
    minimap: { enabled: false },
    scrollBeyondLastLine: false,
    automaticLayout: true,
    wordWrap: "on",
  }}
/>
\end{lstlisting}

Beside the editor, students can run their code two ways. ``Run Code'' hits the
\texttt{/execute} playground with custom stdin and streams back stdout, stderr,
compile output, and time/memory figures --- a scratchpad that persists no
submission. ``Run Samples'' calls \texttt{POST /submissions/run}, which
executes the solution against only the problem's \emph{sample} test cases and
returns per-test pass/fail, so a student gets immediate feedback without ever
seeing the hidden tests that decide their grade. Both paths are guarded by the
same ownership logic as grading: a student may run only against a problem
belonging to an exam assigned to a class they are enrolled in.

\section{Notifications and Announcements}

APEX has no message broker; notifications are plain rows written synchronously
in the request that causes them, with an optional email side-channel. The single
choke point is \texttt{notification.Create}, called from every feature that
needs to reach a user. It inserts a \texttt{Notification} row and then, in a
detached goroutine, checks the recipient's \texttt{notify\_email} preference and
sends an email --- HTML-escaped in both subject-adjacent fields to prevent the
injection class that a pre-ship security review had flagged:

\begin{lstlisting}[language=Go,caption={notification.Create: one in-app row, best-effort escaped email.},label={lst:notif}]
if err := database.DB.Create(&notif).Error; err != nil {
    return
}
go func() { // async email if the user opted in
    var prof models.UserProfile
    if err := database.DB.Where("user_id = ?", userID).First(&prof).Error; err != nil {
        return
    }
    if !prof.NotifyEmail {
        return
    }
    // ... load user, then send with html.EscapeString on title & body
    body := fmt.Sprintf("<p><strong>%s</strong></p><p>%s</p>",
        html.EscapeString(title), html.EscapeString(description))
    _ = email.Send(user.Email, subject, body, text)
}()
\end{lstlisting}

Announcements are the fan-out case. When a teacher posts an announcement to a
class (\texttt{announcement.Create}, gated to the class owner), the handler
persists the announcement row and then loops over every enrolled student,
calling \texttt{notification.Create} once per member with a deep link back to
the class's announcements tab. Each recipient thus gets both an in-app row and,
if opted in, an email --- a genuine fan-out driven by the enrollment table
rather than a broadcast primitive. The same fan-out shape recurs elsewhere: a
student joining a class notifies the teacher, and a graded exam notifies the
student, each through the same \texttt{Create} helper.

On the client, notifications are polled rather than pushed. The
\texttt{useNotificationsQuery} and \texttt{useUnreadCount} hooks are React Query
subscriptions with a 30\,s \texttt{refetchInterval}, enabled only
when a token is present. Marking one or all as read mutates through the API and
invalidates both query keys, so the badge count and the list stay consistent
without a socket connection. Polling is a deliberate simplicity trade: it costs
a lightweight request every half-minute in exchange for needing no persistent
connection infrastructure.

\section{The In-Process Reminder Scheduler}

Exam reminders are the one piece of genuinely time-driven behavior in the
system, and they run inside the same process as the API rather than as a
separate cron service. \texttt{reminder.Start}, invoked once from
\texttt{main.go}, launches a goroutine that sleeps 15\,s to let the
database connection settle and then ticks every minute on a
\texttt{time.Ticker}. Each tick calls \texttt{runOnce}, which is wrapped in a
\texttt{recover} so a panic in one sweep --- a malformed row, a transient query
error --- logs and is contained rather than killing the ticker.

\cref{fig:reminder} lays out the three column-gated sweeps that each tick
performs, positioned relative to an exam's start time~$T$.

\begin{figure}[htbp]\centering
\resizebox{\textwidth}{!}{%
\begin{tikzpicture}[
    font=\small,
    milestone/.style={circle, draw=apexblue, fill=apexlight, thick, minimum size=7pt, inner sep=1.5pt},
    ev/.style={apexbox, align=center, font=\scriptsize, minimum height=9mm, text width=30mm},
    guard/.style={apexstore, align=center, font=\scriptsize, text width=38mm, minimum height=7mm},
    tick/.style={apexgrey, font=\scriptsize},
]

% ---- the number line ----
\draw[apexarrowblue, line width=1.1pt] (-7.6,0) -- (7.6,0);
\node[tick, right] at (7.7,0) {time};
\foreach \x in {-6,0} { \draw[apexblue, thick] (\x,0.14) -- (\x,-0.14); }
\node[tick, below=3pt] at (-6,0) {$T-60\,$min};
\node[tick, below=3pt] at (0,0)  {$T=0$ (exam start)};
\node[milestone] (mA) at (-6,0) {};
\node[milestone] (mB) at (0,0) {};

% ---- events ABOVE the line (kept clear of one another) ----
\node[ev] (inapp) at (-6.6,2.4)
  {\textbf{In-app reminder}\\ \emph{exam\_reminder} fan-out\\ to enrolled students};
\node[guard, above=3mm of inapp] (gA1)
  {guard \texttt{reminder\_sent\_at}\\ null in $(\text{now},\text{now}{+}60\text{m}]$};
\draw[apexarrow] (mA) -- (inapp.south);

\node[ev] (email1h) at (-2.3,2.4)
  {\textbf{1-hour email}\\ subject: \emph{starts in}\\ \emph{about 1 hour}};
\node[guard, above=3mm of email1h] (gA2)
  {guard \texttt{email\_reminder1h\_sent\_at}\\ null in $[\text{now}{+}58\text{m},\text{now}{+}62\text{m}]$};
\draw[apexarrow] (mA) -- (email1h.south);

\node[ev] (start) at (2.2,2.4)
  {\textbf{Start email}\\ subject: \emph{is}\\ \emph{starting now}};
\node[guard, above=3mm of start] (gB)
  {guard \texttt{email\_reminder\_start\_sent\_at}\\ null in $[\text{now}{-}2\text{m},\text{now}{+}2\text{m}]$};
\draw[apexarrow] (mB) -- (start.south);

% ---- the driving ticker BELOW the line (nothing else below to cross) ----
\node[apexext, align=center, font=\scriptsize, text width=62mm] (ticker) at (-1.4,-2.5)
  {\textbf{Scheduler} (\texttt{reminder.Start})\\ 15\,s warm-up, then a 1-minute \texttt{time.Ticker};\\
   each tick runs three column-gated sweeps over\\ published exams (\texttt{is\_draft=false}, \texttt{start\_time} set)};
\draw[apexarrowblue, dashed] (ticker.north west) .. controls (-5.8,-1.7) and (-6.3,-0.9) .. (mA.south);
\draw[apexarrowblue, dashed] (ticker.north east) .. controls (0.9,-1.6) and (0.3,-0.8) .. (mB.south);

% ---- idempotency note (clear right-bottom area) ----
\node[apexlabel, align=center, font=\scriptsize, text width=44mm] at (5.0,-2.5)
  {each timestamp, once non-null, makes its\\ reminder fire \textbf{at most once} --- idempotent\\ across process restarts};

\end{tikzpicture}%
}
\caption{Idempotent exam-reminder schedule relative to start time $T$.}\label{fig:reminder}
\end{figure}


All three sweeps share the same idempotency mechanism, and it is the design's
central idea: each reminder is guarded by a dedicated ``sent-at'' column, and a
reminder fires only for exams where that column is still \texttt{NULL}. After
firing, the column is stamped. Because the guard is a persisted column and not
in-memory state, a process restart cannot cause a re-send --- the reminders are
idempotent across restarts, not merely within a single process lifetime.

\begin{enumerate}
  \item \textbf{In-app, T$-$60.} The main query in \texttt{runOnce} selects
    published exams (\texttt{is\_draft = false}, \texttt{start\_time} not null)
    whose start falls in \texttt{(now, now+60min]} and whose
    \texttt{reminder\_sent\_at} is \texttt{NULL}. For each, \texttt{notifyExam}
    resolves the enrolled students, filters to those with
    \texttt{notify\_exam\_reminders = true}, and fans out an
    \texttt{exam\_reminder} notification linking to \texttt{/dashboard/exams}.
    Even an exam with no classes or members gets its column stamped, so it is
    not re-scanned on every subsequent tick.
  \item \textbf{Email, T$-$60.} \texttt{emailSweep1Hour} selects exams whose
    start falls in the tight \texttt{[now+58min, now+62min]} window with
    \texttt{email\_reminder1h\_sent\_at} still \texttt{NULL}, filters recipients
    by \texttt{notify\_exam\_email}, and sends the ``starts in about 1 hour''
    email before stamping the column.
  \item \textbf{Email, T$-$0.} \texttt{emailSweepStart} does the same for the
    \texttt{[now-2min, now+2min]} window around the start itself, guarded by
    \texttt{email\_reminder\_start\_sent\_at}, with the ``is starting now''
    subject.
\end{enumerate}

The narrow windows on the email sweeps (four minutes wide) combine with the
one-minute tick to guarantee each email is evaluated several times but sent
once. If \texttt{SMTP\_HOST} is unset, \texttt{email.Send} logs the message body
to stdout instead of dispatching it, so the whole reminder path is exercisable
in local development without a mail server. Notably, the ``Exam Graded''
notification is \emph{not} part of this scheduler --- it belongs to the grading
aggregator of \cref{ch:ch07}, which is the only place that knows when every
submission in an attempt has finished grading.

\section{Frontend Architecture}

The frontend is a React~18~\cite{react} single-page application built with
Vite~\cite{vite} and written in TypeScript~\cite{typescript}, totaling roughly
17,400 lines. Its organizing principle is \emph{feature folders}: rather
than grouping by technical kind (all components here, all pages there),
\texttt{src/features/} holds self-contained slices --- \texttt{auth},
\texttt{courses}, \texttt{dashboard}, \texttt{exams}, \texttt{grading},
\texttt{landing}, \texttt{playground}, \texttt{results}, \texttt{settings}, and
\texttt{social} --- each bundling its own pages, components, types, and helpers.
The \texttt{exams} feature seen throughout this chapter is the largest, and its
internal split into \texttt{pages/}, \texttt{components/}, \texttt{lib/}, and
\texttt{types/} is representative.

\subsection{The API Layer and Data Hooks}

All network access funnels through a single \texttt{apiFetch} helper in
\texttt{src/lib/api.ts}. It attaches the bearer token from
\texttt{localStorage["kernel-token"]}, sets the JSON content type, and
centralizes error handling. Its most important behavior is the global
\texttt{401} handler: on an unauthorized response, if a token was present, it is
a stale session --- the helper clears the stored token and role and redirects to
\texttt{/auth}; if no token was present, it simply throws so a public page can
handle it. Every other non-2xx response is unwrapped into a typed
\texttt{ApiError} carrying the backend's error message and status. The rest of
the file is a flat catalogue of typed, one-line functions --- one per endpoint
--- so a caller writes \texttt{getStudentExam(id)} rather than assembling a URL.

On top of that thin layer sits TanStack React Query~\cite{reactquery}, wrapped
in small feature hooks under \texttt{src/hooks/}. This is where the frontend's
server-state discipline lives. Reads are \texttt{useQuery} subscriptions keyed
by a stable tuple --- \texttt{["exam", id]}, \texttt{["student-exams"]},
\texttt{["my-attempts"]} --- and often \texttt{enabled} conditionally on role so
a student never fires a teacher-only query. Writes are \texttt{useMutation}
hooks that invalidate the affected keys on success, letting React Query refetch
and re-render rather than manually reconciling caches:

\begin{lstlisting}[language=TypeScript,caption={A read hook and a write hook: React Query keys plus invalidation.},label={lst:hooks}]
export function useStudentExam(id: number) {
  return useQuery({
    queryKey: ["student-exam", id],
    queryFn: () => api.getStudentExam(id),
    enabled: !!id,
  });
}

export function useUpdateExam() {
  const qc = useQueryClient();
  return useMutation({
    mutationFn: ({ id, data }: { id: number; data: any }) => api.updateExam(id, data),
    onSuccess: (_d, vars) => {
      qc.invalidateQueries({ queryKey: ["exams"] });
      qc.invalidateQueries({ queryKey: ["exam", vars.id] });
    },
  });
}
\end{lstlisting}

The exam-taking flow shows this coordination end to end. When a student first
opens an exam, the component calls \texttt{startExamAttempt} and then
invalidates \texttt{["my-attempts"]} and \texttt{["student-exams"]}, so the
dashboard's ``Start'' button becomes ``Continue'' and the results view learns of
the in-progress attempt --- all without prop-drilling or a global store, because
the query cache is the shared state.

\subsection{UI, Styling, and Routing}

The visual layer is built on shadcn/ui~\cite{shadcn} components --- roughly
thirty Radix-based primitives (dialogs, selects, popovers, switches, scroll
areas, resizable panels) copied into the repository and styled locally rather
than pulled from a versioned package --- over Tailwind CSS~\cite{tailwind}.
Tailwind's utility classes carry essentially all of the styling, including the
conditional emphasis that matters during an exam: the countdown timer, for
instance, switches to a destructive color once fewer than five minutes remain,
expressed inline as a conditional class rather than a separate stylesheet.
Markdown in problem statements and descriptions is rendered with
\texttt{react-markdown} and \texttt{remark-gfm}, deliberately without raw-HTML
passthrough so a problem statement cannot smuggle a script tag.

Routing is handled by \texttt{react-router-dom}~6. The application splits along
role: a public shell serves the landing, authentication, and password-reset
routes, while an authenticated dashboard shell nests the feature pages under
\texttt{/dashboard}. The result is a codebase where a single feature --- exams
--- can be read almost top to bottom: its routes on the backend, its handlers,
its React Query hooks, its pages and editors, and the two-tier persistence that
ties the timed session together, each layer thin enough to hold in one's head
and each boundary drawn where a real responsibility changes hands.
